\lecture{2}{Вопросы на уд}

\section{Вопросы на уд}

\subsection{Классы и базовая ООП механика}
\subsubsection{Модификаторы доступа (public/private)}
\subsubsection{Инициализация по умолчанию}
\subsubsection{Конструкторы и деструкторы}
\subsubsection{Делигирующие конструкторы}
\subsubsection{Списки инициализации}
\subsubsection{Порядок вызова конструкторов и деструкторов}
\subsubsection{Конструктор по умолчанию}
\subsubsection{Конструктор копирования}
\subsubsection{Константные методы}
\subsubsection{Статические поля и методы}
\subsubsection{Оператор присваивания}
\subsubsection{Правило трёх}
\subsubsection{Перегружаемые операторы}
\subsubsection{Перегрузка арифметических операторов и её нюансы}

\subsection{Приведения типов}
\subsubsection{c-style cast}
\subsubsection{static\_cast}
\subsubsection{const\_cast}
\subsubsection{reinterpret\_cast}
\subsubsection{Использование и применимость различных кастов}

\subsection{Шаблоны}
\subsubsection{Основная идея и синтаксис}
\subsubsection{Шаблонные функции}
\subsubsection{Шаблонные классы}
\subsubsection{Псевдонимы, константы и методы}
\subsubsection{Шаблонные методы шаблонного класса}
\subsubsection{Non-type template parameters}

\subsection{Наследование}
\subsubsection{Общая идея}
\subsubsection{Видимость и доступность при публичном наследовании}
\subsubsection{Порядок вызова конструкторов/деструкторов}
\subsubsection{Приведение типов при публичном наследовании}
\subsubsection{Срезка при копировании}

\subsection{Динамический полиморфизм}
\subsubsection{Идея и простой пример}
\subsubsection{Виртуальные функции}
\subsubsection{Абстрактные классы и чисто виртуальные функции}
\subsubsection{Проблема виртуального деструктора}

\subsection{Лямбда-функции}
\subsubsection{Основная идея и синтаксис}
\subsubsection{Примеры}

\subsection{Исключения}
\subsubsection{Идея и примеры использования}
\subsubsection{Разница между исключениями и RE}
\subsubsection{Exception-safety}

\subsection{Итераторы}
\subsubsection{Идея и синтаксис использования}
\subsubsection{Range-based for}
\subsubsection{Виды итераторов и операции}

\subsection{Стандартные контейнеры}
\subsubsection{\texttt{vector}, \texttt{list}, \texttt{deque}, \texttt{unordered\_map}, \texttt{map}, \texttt{set}}
\subsubsection{Основные методы и их асимптотика}
\subsubsection{Адаптеры (на примере \texttt{std::stack})}

\subsection{Move-семантика}
\subsubsection{Проблема и решение}
\subsubsection{Использование \texttt{std::move}}
\subsubsection{Move-конструкторы и оператор перемещающего присваивания (на примере \texttt{std::string})}
\subsubsection{Правило пяти}
\subsubsection{Правильная реализация \texttt{swap}}

\subsection{Умные указатели}
\subsubsection{Идея}
\subsubsection{\texttt{unique\_ptr}, \texttt{shared\_ptr}, \texttt{weak\_ptr}: разница и примеры}

\subsection{Ключевое слово \texttt{auto}} \label{auto}
\subsubsection{Использование для объявления переменных: преимущества и недостатки}
\subsubsection{Правила вывода типа для \texttt{auto}}
\subsubsection{Особенности \texttt{auto\&}, \texttt{const auto\&}}
\subsubsection{Возвращаемый тип функций: trailing return type}
\subsubsection{\texttt{auto} в параметрах функций}
